% !TEX program = xelatex
\documentclass[UTF8,a4paper,12pt,fontset=fandol]{ctexart}

\usepackage{amsmath,amssymb,amsthm,bm}
\usepackage{geometry}
\usepackage{booktabs}
\usepackage{longtable}
\usepackage{enumitem}
\usepackage{xcolor}
\usepackage{hyperref}
\usepackage{url}
\usepackage{algorithm}
\usepackage{algpseudocode}
\usepackage{array}
\usepackage{multirow}

\geometry{left=2.6cm,right=2.6cm,top=2.6cm,bottom=2.6cm}
\hypersetup{
    colorlinks=true,
    linkcolor=blue!60!black,
    citecolor=blue!60!black,
    urlcolor=blue!60!black
}

\newtheorem{definition}{定义}
\newtheorem{assumption}{假设}
\newtheorem{theorem}{定理}
\newtheorem{proposition}{命题}
\newtheorem{lemma}{引理}
\newtheorem{remark}{备注}

\title{\textbf{博士研究开题报告}\\[0.5em]
\Large Language-Guided Verifiable Algorithm Unrolling for Operator-Splitting Optimisers\\
\large 面向组合优化专业求解器的语言引导可验证算法展开}
\author{拟研究方向：大模型辅助优化算法设计、算法展开、组合优化专业求解器}
\date{\today}

\begin{document}
\maketitle

\begin{abstract}
组合优化问题广泛存在于物流调度、能源系统、网络设计、生产计划、广告分配和芯片设计等真实工业场景中。现代专业求解器（如 SCIP、HiGHS、OR-Tools、Gurobi、CPLEX 等）通常依赖分支定界/分支切割、线性或二次松弛、启发式、切平面、预处理与局部搜索等模块协同工作。近年来，大模型在自动算法设计、组合优化启发式生成、程序搜索和反思式进化中显示出潜力；与此同时，将 PDHG、ADMM、proximal gradient 等经典 operator-splitting 优化算法展开为可学习迭代网络，也成为 learning-to-optimise 和 differentiable optimisation 的重要路线。然而，现有研究仍存在明显断层：大模型生成的组合优化算法多偏启发式，缺少与专业求解器内部松弛求解、证书和稳定性的深度结合；而现有 algorithm unrolling 工作多依赖人工设计结构，尚未形成面向组合优化专业求解器的可验证自动设计框架。

本开题报告提出一个博士研究课题：\textbf{Language-Guided Verifiable Algorithm Unrolling for Operator-Splitting Optimisers}。核心思想是：不让大模型直接求解组合优化问题，而让大模型作为离线算法设计器，在受约束的算法领域特定语言（DSL）中生成可验证的展开策略，包括步长调度、罚参数更新、restart 规则、残差平衡、混合算子切换和松弛-修复接口；随后通过静态验证器、理论安全类、执行反馈和多层 benchmark 对候选策略进行筛选。研究目标是将该框架嵌入组合优化专业求解器中的 LP/QP 松弛、拉格朗日松弛、Benders/ADMM 分解子问题和 LNS 子问题，形成兼具可解释性、可验证性、跨规模泛化和工程实用性的 LLM-guided unrolled optimiser。

预期贡献包括：（1）建立面向 operator-splitting optimiser 的语言引导可验证算法搜索空间；（2）提出大模型生成、验证器过滤、执行反馈进化和小规模到大规模迁移的闭环方法；（3）给出安全策略类下的收敛性、残差稳定性、松弛误差传递和求解器集成理论；（4）构建面向组合优化专业求解器的系统 benchmark 与消融协议；（5）在 MIPLIB、Distributional MIPLIB、ML4CO benchmark、CO-Bench/HeuriGym 相关任务及若干真实 MILP/CO 问题上验证方法效果。
\end{abstract}

\tableofcontents
\newpage

\section{研究背景与意义}

\subsection{组合优化专业求解器的现实需求}

组合优化（Combinatorial Optimisation, CO）问题通常可以通过混合整数线性规划（MILP）、混合整数二次规划（MIQP）、约束规划（CP）、网络流、集合覆盖、车辆路径、最大独立集、调度和装箱等形式表达。以 MILP 为例，其标准形式可写为：
\begin{equation}
\begin{aligned}
\min_{x \in \mathbb{R}^{n}} \quad & c^\top x \\
\text{s.t.} \quad & A x \le b,\\
& \ell \le x \le u,\\
& x_i \in \mathbb{Z}, \quad i \in \mathcal{I}.
\end{aligned}
\end{equation}
现代专业求解器通常不会单靠一种算法，而是将多个模块组合为复杂系统：
\begin{itemize}[leftmargin=2em]
    \item \textbf{松弛求解器}：在每个节点求解 LP/QP 松弛，提供 lower bound、dual 信息、reduced cost 和分支信息；
    \item \textbf{分支定界/分支切割}：通过分支、cut separation、domain propagation 和剪枝控制搜索树；
    \item \textbf{启发式与修复}：产生高质量可行解，改善 primal bound；
    \item \textbf{大邻域搜索（LNS）}：固定部分变量并重解子问题；
    \item \textbf{分解与并行化}：对网络流、调度、能源系统等结构化问题使用 ADMM、Benders、Lagrangian relaxation 等。
\end{itemize}

这类专业求解器的优势在于可靠、可解释、具有证书；劣势在于大量策略依赖人工经验，且不同问题族的最优参数和策略差异显著。因此，一个自然问题是：

\begin{quote}
能否利用大模型的程序生成和算法设计能力，自动生成可验证、可集成、可迁移的优化算法展开策略，从而增强专业组合优化求解器？
\end{quote}

\subsection{大模型辅助算法设计的兴起}

近年来，LLM 不再只被用于自然语言任务，也开始被用于程序搜索、数学发现和组合优化启发式生成。FunSearch 通过将预训练大模型与系统化 evaluator 结合，在函数空间中搜索代码形式的解法，强调 evaluator 对幻觉和错误方案的约束作用。EoH 使用大模型和进化计算共同演化启发式思想与可执行代码，面向组合优化启发式自动设计。ReEvo 进一步将反思式语言反馈与进化搜索结合，提出 LLM-based hyper-heuristics。LLM-LNS 将 LLM 用于 MILP 的大邻域搜索策略设计，说明 LLM 设计求解器组件具有现实潜力。CO-Bench 和 HeuriGym 则开始系统评估大模型在组合优化算法搜索与启发式生成上的能力。

这些工作说明：\textbf{LLM 作为算法设计器}是一条正在形成的研究路线。然而，当前 LLM-for-CO 研究主要集中在启发式、元启发式或完整代码生成，常见不足包括：
\begin{enumerate}[leftmargin=2em]
    \item 输出空间过于开放，难以保证稳定性和可验证性；
    \item 与专业求解器内部 LP/QP 松弛、dual bound、certificate、node processing 的结合不足；
    \item 大模型贡献容易退化为昂贵的黑盒超参数搜索；
    \item 缺少跨求解器、跨问题族、跨规模的系统评估协议。
\end{enumerate}

\subsection{算法展开与 operator-splitting optimiser}

Algorithm unrolling/unfolding 将经典迭代算法的有限步展开成可学习结构。对优化算法而言，常见形式为：
\begin{equation}
z_{k+1} = \mathcal{T}_{\theta_k}(z_k; P), \quad k=0,\ldots,K-1,
\end{equation}
其中 $P$ 是问题实例，$z_k$ 是算法状态，$\mathcal{T}_{\theta_k}$ 是带参数的更新算子。与纯黑盒神经网络不同，algorithm unrolling 具有结构清晰、可解释、易结合理论和可控制迭代步数等优点。

在组合优化专业求解器中，最适合展开的不是 simplex 或 branch-and-bound 主循环本身，而是以下 operator-splitting 模块：
\begin{itemize}[leftmargin=2em]
    \item \textbf{PDHG / PDLP}：适合大规模 LP 松弛，矩阵向量乘友好，可部署于 GPU；
    \item \textbf{ADMM / Douglas--Rachford splitting}：适合可分结构、分解式约束和 Lagrangian relaxation；
    \item \textbf{Proximal gradient / FISTA}：适合正则化松弛、二次罚函数和连续 surrogate；
    \item \textbf{Augmented Lagrangian / SQP 子问题}：适合约束优化与 MIQP/NLP 松弛；
    \item \textbf{PCG/linear solver module}：适合 QP、Newton、Gauss--Newton 或 KKT reduced system 的内层线性求解。
\end{itemize}

这些算法天然拥有层状结构、残差信息、步长/罚参数/restart 规则和可验证的安全条件。因此，相比直接让 LLM 生成任意启发式，\textbf{在 operator-splitting optimiser 的受控搜索空间内使用 LLM 自动设计展开策略}更有希望形成兼具新意和可靠性的 AI 会议论文方向。

\section{相关工作综述}

\subsection{LLM 自动算法设计与组合优化}

\paragraph{FunSearch.}
FunSearch 将 LLM 与自动 evaluator 配对，在代码函数空间中进行演化搜索，强调“LLM 提出候选，evaluator 负责验证”的闭环范式。该思想对本课题具有基础性启发：LLM 不应直接被当作 solver，而应作为 candidate generator。

\paragraph{EoH 与 ReEvo.}
EoH 将启发式的自然语言思想与可执行代码同时进化；ReEvo 通过反思式语言反馈进一步提升 heuristic generation 的效率和质量。二者共同说明，LLM 可以利用已有算法知识生成有价值的组合优化启发式。然而，这类方法通常输出开放式 heuristic，缺少与 solver certificate 和 operator-splitting 更新条件的紧密绑定。

\paragraph{LLM-LNS.}
LLM-LNS 将 LLM 用于 MILP 中的大邻域搜索策略设计，与 Gurobi、SCIP、手工 LNS 和学习型 LNS 进行比较。该工作最接近本课题的“专业求解器组件设计”动机，但其核心仍偏 neighborhood selection heuristic，而不是面向 LP/QP 松弛或 operator-splitting unrolled optimiser 的可验证更新设计。

\paragraph{CO-Bench 与 HeuriGym.}
CO-Bench 和 HeuriGym 说明研究社区正在重视 LLM 在组合优化算法搜索中的系统评估。这对本课题的 benchmark 设计具有直接参考价值：必须设置明确目标函数、可执行反馈、质量指标、失败率和样本效率指标，避免只做 narrative demo。

\subsection{Algorithm unrolling 与 learned optimisation}

\paragraph{PDHG-Net.}
PDHG-Unrolled Learning-to-Optimize Method for Large-Scale LP 将 PDHG 展开为 PDHG-Net，并采用“两阶段”方法：先由网络生成近似解，再由 PDHG 精修。其重要启发是：AI 会议认可“经典一阶优化器 + 展开结构 + 可靠后处理”的混合求解模式。

\paragraph{DeepDistributedQP.}
DeepDistributedQP 将 OSQP/consensus 风格的 QP 求解器展开成深度架构，并通过 PAC-Bayes 理论给出泛化边界。该工作表明：operator-splitting optimiser 的 learned unfolding 可以与大规模 QP、跨规模泛化和理论保证结合。

\paragraph{Deep FlexQP.}
Deep FlexQP 通过 deep unfolding 加速 nonlinear programming 中的 QP 子问题，并强调可行性、松弛与求解器集成。该方向说明，unrolled solver 若想走向真实专业求解器，必须关注 infeasibility、constraint violation 和 solver-compatible certificates，而不仅是 objective loss。

\subsection{ML for CO 与专业求解器增强}

ML4CO 研究包括 branching policy learning、cut selection、node selection、primal heuristic、LNS、neural diving、neural rounding、solver configuration 和 instance embedding 等。已有方法多使用 GNN 表示变量--约束二部图，并学习求解器局部决策。与之相比，本课题更关注：
\begin{itemize}[leftmargin=2em]
    \item 不直接替代 solver 的关键决策，而是设计可验证的松弛求解和分解优化模块；
    \item 不只学习数值参数，而是通过 LLM 生成结构化更新规则；
    \item 不追求端到端黑盒，而是保留 operator-splitting 的理论骨架；
    \item 强调与专业 solver 的安全接口：bound、feasibility、certificate、repair。
\end{itemize}

\section{问题陈述与研究目标}

\subsection{核心研究问题}

本课题拟研究如下核心问题：

\begin{quote}
如何利用大语言模型，在可验证的受控算法空间中自动设计 operator-splitting optimiser 的展开策略，并将其嵌入组合优化专业求解器，以提升松弛求解、启发式修复、分解子问题和 LNS 子问题的效率与泛化能力？
\end{quote}

该问题可分解为四个研究问题：

\begin{description}[leftmargin=2em]
    \item[RQ1：搜索空间] 如何定义适合 LLM 生成、可执行、可验证、可理论分析的 unrolled optimiser DSL？
    \item[RQ2：安全验证] 如何保证 LLM 生成的策略不破坏基本优化性质，如有界步长、残差稳定、proximal non-expansiveness、feasibility repair 和 bound validity？
    \item[RQ3：求解器集成] 如何将生成策略嵌入 MILP/CO solver 的 LP/QP 松弛、Lagrangian/ADMM 分解、LNS 子问题和 rounding/repair 模块？
    \item[RQ4：实验评估] 如何设计 benchmark 和 baseline，证明方法不是 prompt engineering、不是黑盒超参数搜索、不是只在 toy problem 上有效？
\end{description}

\subsection{研究对象}

本文将组合优化专业求解器中的候选模块抽象为：
\begin{equation}
\mathcal{S}(P;\pi) = \texttt{SolverCore}\Big(P, \texttt{RelaxSolver}_{\pi}, \texttt{Heuristics}_{\pi}, \texttt{Repair}, \texttt{Certify}\Big),
\end{equation}
其中 $P$ 为组合优化实例，$\pi$ 是由 LLM 生成并通过 verifier 筛选的策略程序。$\pi$ 不直接输出解，而是输出以下组件之一或多个：
\begin{itemize}[leftmargin=2em]
    \item PDHG/ADMM/FISTA 的步长、罚参数、momentum 和 restart 调度；
    \item 残差平衡规则和 adaptive scaling；
    \item operator switching 规则，例如从 PDHG 切换到 proximal refinement；
    \item LNS 中 destroy/repair 的 schedule，与松弛残差、reduced cost 或 incumbent improvement 绑定；
    \item rounding-aware penalty update 与 feasibility repair 接口；
    \item 可选的 PCG/linear solver restart、truncation 或 preconditioning policy。
\end{itemize}

\subsection{总体目标}

本课题的总体目标是提出一个可复用框架：
\begin{equation}
\boxed{
\text{LLM} + \text{DSL} + \text{Verifier} + \text{Evaluator} + \text{Solver Integration}
}
\end{equation}
使得大模型生成的算法展开策略能够：
\begin{enumerate}[leftmargin=2em]
    \item 在小规模实例上通过自动评估和反思进化快速筛选；
    \item 在大规模和分布外实例上保持稳定；
    \item 在专业求解器中提供可解释的加速或质量提升；
    \item 在理论上满足某类安全条件；
    \item 在评估上战胜手工规则、随机搜索、贝叶斯优化、进化搜索、常规 learned schedule 和 LLM heuristic generation baseline。
\end{enumerate}

\section{拟提出的方法框架}

\subsection{总体框架}

整体流程如算法 \ref{alg:framework} 所示。

\begin{algorithm}[h]
\caption{Language-Guided Verifiable Algorithm Unrolling}
\label{alg:framework}
\begin{algorithmic}[1]
\Require Problem family $\mathcal{D}$, algorithm templates $\mathcal{A}$, DSL grammar $\mathcal{G}$, verifier $\mathcal{V}$, evaluator $\mathcal{E}$, LLM $\mathcal{M}$
\State Initialise candidate pool $\mathcal{P}_0$ with human-safe baseline policies
\For{$t=0,1,\ldots,T-1$}
    \State Select elite policies from $\mathcal{P}_t$ according to validation score
    \State Prompt LLM $\mathcal{M}$ with problem description, DSL grammar, failure cases, and elite policies
    \State Generate candidate policies $\{\pi_i\}_{i=1}^{m} \subset \mathcal{G}$
    \For{each $\pi_i$}
        \State Apply static verifier: type check, bound check, operator legality
        \State Apply mathematical verifier: safe parameter ranges, monotonicity proxy, proximal validity
        \State Apply smoke-test evaluator on small instances
        \If{$\mathcal{V}(\pi_i)=\texttt{pass}$}
            \State Evaluate $\pi_i$ on training instances using $\mathcal{E}$
            \State Add $\pi_i$ to candidate pool
        \EndIf
    \EndFor
    \State Construct reflection messages from failures and successful cases
    \State $\mathcal{P}_{t+1} \leftarrow$ elites + diverse verified candidates
\EndFor
\State Return best policies and optionally fine-tuned scalar constants
\end{algorithmic}
\end{algorithm}

框架关键点如下：
\begin{itemize}[leftmargin=2em]
    \item \textbf{LLM 离线设计}：LLM 只在策略搜索阶段调用，部署求解时不调用 LLM；
    \item \textbf{受控 DSL}：LLM 只能生成有限、安全、可解释的策略程序；
    \item \textbf{多级 verifier}：先过滤不可运行和不安全策略，再进行 expensive benchmark；
    \item \textbf{执行反馈闭环}：候选策略通过真实优化指标排序，而不是通过语言判断；
    \item \textbf{求解器兼容}：最终策略必须能够嵌入专业 solver 的 relaxation、heuristic 或 LNS 模块。
\end{itemize}

\subsection{Operator-splitting optimiser 模板}

\subsubsection{PDHG/PDLP 模板}

考虑 saddle-point 形式：
\begin{equation}
\min_{x \in X} \max_{y \in Y} \; f(x) + \langle Kx, y\rangle - g^\ast(y).
\end{equation}
标准 PDHG 更新为：
\begin{equation}
\begin{aligned}
y^{k+1} &= \operatorname{prox}_{\sigma_k g^\ast}\left(y^k + \sigma_k K\bar{x}^k\right),\\
x^{k+1} &= \operatorname{prox}_{\tau_k f}\left(x^k - \tau_k K^\top y^{k+1}\right),\\
\bar{x}^{k+1} &= x^{k+1} + \theta_k(x^{k+1}-x^k).
\end{aligned}
\end{equation}
LLM 可设计：
\begin{equation}
(\tau_k,\sigma_k,\theta_k,\texttt{restart}_k,\texttt{rescale}_k)
= \pi_{\text{PDHG}}(s_k),
\end{equation}
其中 $s_k$ 包括 primal residual、dual residual、duality gap proxy、objective change、iteration index、constraint violation 和 bound progress。

\subsubsection{ADMM 模板}

考虑可分问题：
\begin{equation}
\min_{x,z} f(x)+g(z) \quad \text{s.t.} \quad Bx+Cz=d.
\end{equation}
ADMM 更新为：
\begin{equation}
\begin{aligned}
x^{k+1} &= \arg\min_x f(x) + \frac{\rho_k}{2}\|Bx+Cz^k-d+u^k\|^2,\\
z^{k+1} &= \arg\min_z g(z) + \frac{\rho_k}{2}\|Bx^{k+1}+Cz-d+u^k\|^2,\\
u^{k+1} &= u^k + Bx^{k+1}+Cz^{k+1}-d.
\end{aligned}
\end{equation}
LLM 可生成 residual balancing 规则：
\begin{equation}
\rho_{k+1}=h_\pi(\rho_k,\|r^k_p\|,\|r^k_d\|,\Delta f_k,k),
\end{equation}
并由 verifier 保证：
\begin{equation}
0<\rho_{\min}\le \rho_k \le \rho_{\max}<\infty.
\end{equation}

\subsubsection{Proximal gradient / FISTA 模板}

对连续松弛或正则化 surrogate：
\begin{equation}
\min_x F(x)=\phi(x)+\lambda R(x),
\end{equation}
其中 $\phi$ 光滑，$R$ 可 prox。基本更新为：
\begin{equation}
x^{k+1}=\operatorname{prox}_{\eta_k\lambda R}(v^k-\eta_k \nabla \phi(v^k)).
\end{equation}
LLM 可设计：
\begin{equation}
(\eta_k,\beta_k,\texttt{restart}_k,\texttt{rounding\_trigger}_k)
=\pi_{\text{prox}}(s_k).
\end{equation}

\subsection{算法 DSL 设计}

DSL 的目标是限制 LLM 的输出空间，使其既有表达能力，又能被验证。一个策略程序可写为：
\begin{verbatim}
policy PDHG_Schedule:
    features:
        rp = norm(primal_residual)
        rd = norm(dual_residual)
        gap = gap_proxy
        prog = previous_gap - gap
    update:
        if rp > c1 * rd:
            tau = clip(tau * a1, tau_min, tau_max)
            sigma = clip(sigma / a1, sigma_min, sigma_max)
        elif rd > c2 * rp:
            tau = clip(tau / a2, tau_min, tau_max)
            sigma = clip(sigma * a2, sigma_min, sigma_max)
        if prog < eps and k > k_min:
            restart_momentum()
        enforce tau * sigma * K_norm_sq <= 1 - margin
\end{verbatim}

DSL 包含：
\begin{itemize}[leftmargin=2em]
    \item \textbf{状态特征}：残差、gap proxy、bound progress、incumbent improvement、node depth、constraint violation；
    \item \textbf{安全算子}：clip、normalize、restart、damping、rescale、prox、projection；
    \item \textbf{模板变量}：$\tau,\sigma,\rho,\theta,\eta,\beta,K_{\text{inner}}$；
    \item \textbf{约束语句}：boundedness、step-size product、monotonic proxy、repair requirement；
    \item \textbf{求解器接口}：return relaxed solution, dual certificate, rounding hint, LNS neighbourhood score。
\end{itemize}

\subsection{Verifier 设计}

Verifier 分为四层。

\paragraph{第一层：语法与类型验证。}
检查 DSL 是否可解析，变量是否定义，算子是否在允许集合中，输出是否符合求解器接口。

\paragraph{第二层：静态安全验证。}
检查参数范围、步长积、罚参数上下界、restart 规则是否可能导致非法状态。例如 PDHG 安全条件：
\begin{equation}
\tau_k\sigma_k\|K\|^2 \le 1-\delta.
\end{equation}

\paragraph{第三层：动态 smoke test。}
在极小实例上运行若干步，检查是否出现 NaN、Inf、残差爆炸、不可恢复 infeasibility、时间超限或违反 solver API。

\paragraph{第四层：数学安全类验证。}
将候选策略归入某个安全类 $\Pi_{\text{safe}}$。例如：
\begin{equation}
\Pi_{\text{PDHG}} =
\left\{
\pi:\;
\tau_k,\sigma_k>0,\;
\tau_k\sigma_k\|K\|^2\le 1-\delta,\;
\theta_k\in[0,1],\;
\forall k
\right\}.
\end{equation}
后续理论只需对 $\Pi_{\text{safe}}$ 成立，而不必证明 LLM 本身可靠。

\subsection{Evaluator 设计}

候选策略评分函数应同时考虑质量、速度、稳定性和泛化。对单个实例 $P$，可定义：
\begin{equation}
\operatorname{Score}(\pi;P)
=
w_1 \cdot \operatorname{Gap}(P,\pi)
+
w_2 \cdot \operatorname{Viol}(P,\pi)
+
w_3 \cdot \operatorname{Time}(P,\pi)
+
w_4 \cdot \operatorname{Fail}(P,\pi)
+
w_5 \cdot \operatorname{Cost}_{\text{LLM}}(\pi),
\end{equation}
其中：
\begin{itemize}[leftmargin=2em]
    \item $\operatorname{Gap}$：objective gap、duality gap、integrality gap 或 primal integral；
    \item $\operatorname{Viol}$：约束违反、整数违反、可行性修复失败率；
    \item $\operatorname{Time}$：wall-clock、迭代数、node time；
    \item $\operatorname{Fail}$：数值失败、超时、不可解析、solver crash；
    \item $\operatorname{Cost}_{\text{LLM}}$：LLM 查询成本与候选生成次数。
\end{itemize}

对训练集 $\mathcal{D}_{\text{train}}$，总体评分为：
\begin{equation}
\operatorname{Score}(\pi)
=
\mathbb{E}_{P\sim \mathcal{D}_{\text{train}}}
\left[\operatorname{Score}(\pi;P)\right]
+
\lambda \operatorname{Complexity}(\pi),
\end{equation}
其中 $\operatorname{Complexity}(\pi)$ 惩罚过长、过深、不可解释的策略程序。

\section{面向组合优化专业求解器的集成方案}

\subsection{集成层次}

本课题不试图重写完整商业求解器，而采用由浅入深的三层集成路线。

\subsubsection{Level 1：外部 relaxation accelerator}

将 MILP/CO 问题的 LP/QP 松弛导出为标准格式，用 LLM-generated PDHG/ADMM/FISTA 策略加速松弛求解，输出：
\begin{itemize}[leftmargin=2em]
    \item relaxation objective bound；
    \item approximate primal/dual solution；
    \item reduced-cost-like score；
    \item rounding 或 branching hint。
\end{itemize}
该层最容易实现，适合早期验证。

\subsubsection{Level 2：求解器 callback / primal heuristic}

将生成策略嵌入 SCIP/OR-Tools/HiGHS 等开源求解器的 callback：
\begin{itemize}[leftmargin=2em]
    \item 使用松弛解指导 rounding；
    \item 使用 residual/gap 信息触发 local repair；
    \item 使用 LNS schedule 选择需要释放的变量集合；
    \item 将候选解交给专业 solver repair，保证最终可行性。
\end{itemize}

\subsubsection{Level 3：分解式子问题求解}

对网络流、能源调度、设施选址、多商品流、车辆路径等结构化问题，使用 ADMM/Benders/Lagrangian relaxation 构建分解求解器。LLM 设计：
\begin{itemize}[leftmargin=2em]
    \item penalty update；
    \item block update order；
    \item restart/damping；
    \item subproblem time allocation；
    \item primal recovery 和 repair schedule。
\end{itemize}

\subsection{与分支定界的安全接口}

在专业 MILP solver 中，任何学习模块都不能破坏 soundness。因而本课题采用如下原则：
\begin{enumerate}[leftmargin=2em]
    \item \textbf{不替代 pruning certificate}：学习模块给出的 lower bound 若不满足严格认证，则只能作为启发式信号，不能用于剪枝；
    \item \textbf{可行解必须由 solver 验证}：rounding 或 LNS 产生的解必须经过原始约束检查；
    \item \textbf{松弛误差有记录}：approximate relaxation bound 必须携带 residual/gap certificate；
    \item \textbf{失败可回退}：任何 LLM-generated policy 失败时立即回退到 vanilla PDHG/ADMM/solver default。
\end{enumerate}

\section{理论研究计划}

本课题的理论目标不是证明“大模型会生成好算法”，而是证明：\textbf{只要大模型生成的策略属于某个可验证安全类，则展开优化器保留基本性质；其近似误差在专业求解器中以可控方式传播。}

\subsection{安全类下的 PDHG 稳定性}

\begin{assumption}[凸松弛与有界算子]
考虑凸-凹 saddle-point 问题，$f$ 和 $g^\ast$ 为闭真凸函数，$K$ 为有界线性算子。
\end{assumption}

\begin{definition}[PDHG 安全策略类]
称策略 $\pi$ 属于 $\Pi_{\mathrm{PDHG}}(\delta)$，若其生成的参数满足：
\begin{equation}
\tau_k>0,\quad \sigma_k>0,\quad \theta_k\in[0,1],
\quad
\tau_k\sigma_k\|K\|^2 \le 1-\delta,
\quad \forall k.
\end{equation}
\end{definition}

\begin{theorem}[安全类下的 ergodic gap 收敛目标]
若 $\pi\in\Pi_{\mathrm{PDHG}}(\delta)$，则在标准 PDHG 假设下，展开迭代的 ergodic primal-dual gap 可由 $O(1/K)$ 控制；常数依赖于 $\delta$、初始距离和参数上下界。
\end{theorem}

\begin{remark}
该定理可作为第一阶段理论目标。实际论文中需处理 adaptive $\tau_k,\sigma_k$ 和 restart 对 Lyapunov 函数的影响。核心思想是 verifier 将 LLM 生成策略限制在不破坏 PDHG 基本收敛条件的区域内。
\end{remark}

\subsection{ADMM 罚参数策略的收敛稳定性}

\begin{definition}[ADMM bounded penalty policy]
称 $\pi$ 属于 $\Pi_{\mathrm{ADMM}}(\rho_{\min},\rho_{\max})$，若：
\begin{equation}
0 < \rho_{\min}\le \rho_k \le \rho_{\max}<\infty,
\quad \forall k,
\end{equation}
并且每次 $\rho$ 变化时对 scaled dual variable 进行一致重标定。
\end{definition}

\begin{theorem}[变罚参数 ADMM 的残差稳定性目标]
在闭真凸、鞍点存在和子问题精确求解的条件下，若 $\pi\in\Pi_{\mathrm{ADMM}}$ 且参数变化总变差受控，则 primal residual 和 dual residual 收敛到零，objective value 收敛到最优值。
\end{theorem}

该方向将为 LLM 设计 residual balancing rule 提供理论安全边界。

\subsection{松弛误差向 MILP 求解器的传播}

设某节点 LP 松弛最优值为 $v^\star$，approximate relaxation solver 输出 $\hat v$ 与 gap certificate $\epsilon$，满足：
\begin{equation}
|\hat v - v^\star| \le \epsilon.
\end{equation}

\begin{proposition}[近似松弛 bound 的安全使用]
若当前 incumbent 为 $U$，则只有当
\begin{equation}
\hat v - \epsilon \ge U
\end{equation}
时，才可安全剪枝；若无法给出严格 $\epsilon$，则 $\hat v$ 只能用于 node scoring、branching hint 或 heuristic，而不能用于 correctness-critical pruning。
\end{proposition}

该命题体现专业求解器集成的基本原则：学习模块可以增强效率，但不能破坏求解器正确性。

\subsection{展开策略的跨规模泛化分析}

对于策略程序 $\pi$，其复杂度可以用 DSL 长度、条件分支数、连续参数数目和算子组合深度度量。可考虑 PAC-Bayes 或 algorithmic stability 风格的泛化界：
\begin{equation}
\mathbb{E}_{P\sim\mathcal{D}}\left[\ell(\pi;P)\right]
\le
\frac{1}{N}\sum_{i=1}^N \ell(\pi;P_i)
+
\mathcal{O}\left(
\sqrt{\frac{\operatorname{Complexity}(\pi)+\log(1/\delta)}{N}}
\right).
\end{equation}
该方向可解释为什么受控 DSL 比开放式代码生成更容易泛化。

\section{实验设计}

\subsection{实验原则}

为避免“纯讲故事”，实验必须遵循以下原则：
\begin{enumerate}[leftmargin=2em]
    \item \textbf{固定查询预算}：LLM、进化搜索、随机搜索、BO 使用相同候选评估预算；
    \item \textbf{训练/验证/测试分离}：小规模实例用于搜索，大规模和分布外实例用于最终评估；
    \item \textbf{报告失败率}：包括不可运行、验证失败、数值爆炸、超时和可行性失败；
    \item \textbf{报告求解器指标}：不仅报告 loss，还报告 gap、time、node count、primal integral、feasibility 和 bound；
    \item \textbf{消融 LLM 的贡献}：比较 LLM、随机生成、进化无 LLM、BO、手工规则和 learned controller。
\end{enumerate}

\subsection{Benchmark 选择}

\begin{longtable}{p{3.2cm}p{4.6cm}p{5.2cm}}
\toprule
阶段 & 数据/问题 & 目的 \\
\midrule
Stage 1 & 合成 LP/QP、LASSO、network flow、assignment relaxation & 快速验证 DSL、verifier、PDHG/ADMM/FISTA 策略是否有效 \\
Stage 2 & MIPLIB 2017、Distributional MIPLIB、ML4CO competition-style MILP distributions & 测试与专业 MILP solver 相关的 relaxation、rounding、LNS、node-level 指标 \\
Stage 3 & set cover、facility location、knapsack、MIS、TSP/VRP MILP formulation & 测试结构化 CO 问题上的泛化和 solver integration \\
Stage 4 & CO-Bench、HeuriGym 相关任务 & 与 LLM heuristic generation benchmark 对齐，展示本方法不是普通 LLM heuristic \\
Stage 5 & 真实应用案例：物流、能源调度、广告分配或网络设计 & 检验工程价值与鲁棒性 \\
\bottomrule
\end{longtable}

\subsection{Baseline 设计}

\begin{longtable}{p{3.5cm}p{9.2cm}}
\toprule
Baseline 类型 & 具体方法 \\
\midrule
专业求解器 & Gurobi、CPLEX、SCIP、HiGHS、OR-Tools PDLP/PDHG、OSQP \\
原始算法 & vanilla PDHG、vanilla ADMM、FISTA、Douglas--Rachford、solver default LNS \\
手工自适应规则 & residual balancing ADMM、adaptive restart FISTA、hand-crafted PDHG restart/scaling \\
黑盒搜索 & random search、grid search、Bayesian optimisation、CMA-ES、evolutionary search without LLM \\
学习型 unrolling & learned scalar schedule、MLP/LSTM controller、hypernetwork-generated parameters、GNN policy \\
LLM 算法设计 & FunSearch-style program search、EoH-style heuristic evolution、ReEvo-style reflective evolution、LLM-LNS-style neighbourhood design \\
消融版本 & no verifier、no reflection、no evolution、no DSL constraints、LLM-only、small-scale only、without solver repair \\
\bottomrule
\end{longtable}

\subsection{评价指标}

\paragraph{松弛求解指标}
\begin{itemize}[leftmargin=2em]
    \item primal residual、dual residual、duality gap；
    \item time-to-$\epsilon$ gap；
    \item fixed-budget final gap；
    \item GPU/CPU wall-clock；
    \item memory footprint。
\end{itemize}

\paragraph{MILP/CO 求解器指标}
\begin{itemize}[leftmargin=2em]
    \item primal bound、dual bound、integrality gap；
    \item primal integral；
    \item number of explored nodes；
    \item time to first feasible solution；
    \item time to best incumbent；
    \item final optimality gap under time limit；
    \item feasibility violation 和 repair success rate。
\end{itemize}

\paragraph{LLM algorithm design 指标}
\begin{itemize}[leftmargin=2em]
    \item valid program rate；
    \item verified safe program rate；
    \item query budget；
    \item candidate evaluation budget；
    \item best score vs. number of generated candidates；
    \item transfer performance from small to large instances；
    \item diversity and interpretability of discovered policies。
\end{itemize}

\subsection{关键实验}

\subsubsection{实验一：LLM 是否优于黑盒搜索？}

在固定候选数量 $B$ 下比较：
\begin{equation}
\text{LLM+Verifier+Evolution}
\quad \text{vs.} \quad
\text{Random DSL Search}
\quad \text{vs.} \quad
\text{BO/CMA-ES}
\quad \text{vs.} \quad
\text{Human Rule}.
\end{equation}
核心指标为 best verified score、valid rate 和大规模迁移表现。

\subsubsection{实验二：可验证约束是否必要？}

比较：
\begin{itemize}[leftmargin=2em]
    \item full method；
    \item no static verifier；
    \item no mathematical safety constraints；
    \item open-ended code generation；
    \item LLM natural language prompt only。
\end{itemize}
预期验证器能显著降低 crash rate 和 OOD failure rate。

\subsubsection{实验三：求解器集成收益}

在 MIPLIB/Distributional MIPLIB 子集上测试：
\begin{itemize}[leftmargin=2em]
    \item default solver；
    \item solver + vanilla relaxation accelerator；
    \item solver + hand-crafted unrolled strategy；
    \item solver + LLM-generated verified strategy。
\end{itemize}
指标包括 time-to-incumbent、gap、node count、primal integral。

\subsubsection{实验四：跨规模泛化}

在小规模实例上搜索策略，在大规模实例上测试：
\begin{equation}
n_{\text{train}} \ll n_{\text{test}}, \quad m_{\text{train}} \ll m_{\text{test}}.
\end{equation}
分析策略是否依赖绝对维度，或只依赖归一化残差和比例特征。

\subsubsection{实验五：策略可解释性与失败分析}

对最优策略进行人工解释：
\begin{itemize}[leftmargin=2em]
    \item 它是否类似已知 residual balancing 或 restart 规则？
    \item 它在哪些实例族上失败？
    \item 失败是否由 ill-conditioning、infeasibility、弱松弛、整数结构或 solver callback overhead 引起？
\end{itemize}

\section{预期创新点}

\subsection{从 LLM heuristic generation 到 verifiable unrolled optimiser design}

现有 LLM-for-CO 多生成启发式或完整代码。本课题强调生成的是 operator-splitting optimiser 的结构化展开策略，并由 verifier 和理论安全类约束。这使得方法更适合专业求解器，而不是只做开放式 heuristic search。

\subsection{面向专业求解器的安全接口}

本课题明确区分：
\begin{itemize}[leftmargin=2em]
    \item correctness-critical 信息：必须由严格 solver 或 certificate 保证；
    \item heuristic information：可由学习模块提供，用于 node scoring、branching hint、rounding、LNS；
    \item approximate relaxation 信息：必须携带 residual/gap 证书。
\end{itemize}
这使得 LLM-generated optimiser 能安全嵌入专业 solver。

\subsection{可验证 DSL 与理论安全类}

通过 DSL 控制表达空间，并将候选策略映射到 $\Pi_{\text{safe}}$，避免开放式代码生成的不可控性。理论上证明的是安全类性质，而不是 LLM 本身。

\subsection{跨算法族统一视角}

将 PDHG、ADMM、FISTA、Douglas--Rachford、PCG 等算法统一为：
\begin{equation}
z_{k+1}=\mathcal{T}_{a_k,\theta_k}(z_k;P),
\end{equation}
其中 LLM 生成的是 $a_k$ 和 $\theta_k$ 的策略程序。这比单一算法 unrolling 更具有博士研究容量。

\subsection{组合优化 benchmark 与 LLM cost-aware 评估}

本课题将专业求解器指标与 LLM algorithm design 指标结合，报告 valid rate、query cost、candidate budget、OOD failure 等，使得可行性论证更扎实。

\section{可行性分析}

\subsection{技术可行性}

\begin{itemize}[leftmargin=2em]
    \item PDHG/ADMM/FISTA 等算法已有成熟理论和开源实现，适合作为第一阶段模板；
    \item SCIP、HiGHS、OR-Tools、Pyomo、JuMP 等生态支持求解器集成和 benchmark；
    \item LLM program search 可以先基于 API 调用和离线执行，不需要训练大模型；
    \item DSL 限制输出空间，降低 hallucination 和不可运行风险；
    \item benchmark 可从合成实例逐步扩展到 MIPLIB 和 Distributional MIPLIB。
\end{itemize}

\subsection{研究风险与应对}

\begin{longtable}{p{4cm}p{8.5cm}}
\toprule
风险 & 应对方案 \\
\midrule
LLM 生成策略不优于随机搜索 & 设置相同预算下的 random/BO/evolution baseline；强调 LLM 是否提升 sample efficiency 和 valid rate \\
开放式代码不可控 & 使用 DSL、static verifier 和安全类理论；禁止任意 Python 代码进入 solver core \\
求解器集成过难 & 先做 external relaxation accelerator，再做 callback/LNS，最后考虑深度集成 \\
理论只能覆盖凸松弛 & 明确理论对象为 LP/QP/convex surrogate；对整数层只证明安全接口和误差传播 \\
效果只在小规模有效 & 从一开始设计 train-small-test-large 和 OOD benchmark \\
LLM cost 过高 & 离线生成，部署不调用 LLM；报告 query budget 和 amortised cost \\
审稿人认为是 prompt engineering & 公开 DSL、evaluator、query logs、候选池、失败案例和固定预算协议 \\
\bottomrule
\end{longtable}

\section{研究计划与时间安排}

\subsection{第一阶段：文献、DSL 与基础原型（第 1 年）}

\begin{itemize}[leftmargin=2em]
    \item 系统整理 LLM-for-CO、algorithm unrolling、operator-splitting、ML4CO solver enhancement 文献；
    \item 实现 PDHG、ADMM、FISTA 三个基础模板；
    \item 设计第一版 DSL 和 static verifier；
    \item 在 LASSO、network flow、assignment relaxation 和合成 QP 上验证；
    \item 目标产出：一篇 workshop/CCF-B 论文或技术报告。
\end{itemize}

\subsection{第二阶段：LLM 搜索闭环与理论安全类（第 2 年）}

\begin{itemize}[leftmargin=2em]
    \item 实现 LLM + reflection + evolutionary search；
    \item 完成 dynamic verifier 和 evaluator；
    \item 形成 PDHG/ADMM 安全策略类理论；
    \item 与 random search、BO、evolution、learned controller 系统对比；
    \item 目标产出：AAAI/IJCAI/ICLR 投稿。
\end{itemize}

\subsection{第三阶段：专业求解器集成与组合优化 benchmark（第 3 年）}

\begin{itemize}[leftmargin=2em]
    \item 集成 SCIP/HiGHS/OR-Tools 或 Gurobi callback；
    \item 在 MIPLIB、Distributional MIPLIB、ML4CO-style distributions 上测试；
    \item 设计 LNS/rounding/repair 接口；
    \item 完成跨规模、跨问题族泛化实验；
    \item 目标产出：ICLR/ICML/NeurIPS 或高水平 OR/AI 交叉会议论文。
\end{itemize}

\subsection{第四阶段：统一理论、真实应用与博士论文（第 4 年）}

\begin{itemize}[leftmargin=2em]
    \item 完善泛化理论、误差传播理论和 solver soundness 接口；
    \item 在真实工业问题上验证；
    \item 开源框架与 benchmark 协议；
    \item 汇总形成博士论文。
\end{itemize}

\section{拟投稿方向}

\begin{itemize}[leftmargin=2em]
    \item \textbf{ICLR}：最适合 language-guided unrolled optimiser、结构化学习求解器、跨规模泛化；
    \item \textbf{ICML/NeurIPS}：适合理论更强的 safe policy class、泛化、differentiable optimisation 和算法设计；
    \item \textbf{AAAI/IJCAI}：适合 solver-integrated LLM algorithm design 和组合优化应用；
    \item \textbf{CPAIOR/INFORMS/Mathematical Programming Computation}：适合更偏专业求解器和运筹优化的版本；
    \item \textbf{NeurIPS/ICLR workshop}：适合早期 benchmark、DSL 和系统原型。
\end{itemize}

\section{预期成果}

\begin{enumerate}[leftmargin=2em]
    \item 一个可复用的 \textbf{Language-Guided Verifiable Unrolling} 框架；
    \item 一个面向 operator-splitting optimiser 的 DSL 和 verifier；
    \item 若干被 LLM 发现的可解释 PDHG/ADMM/FISTA/LNS 策略；
    \item 安全策略类下的收敛性、稳定性和误差传播理论；
    \item 面向组合优化专业求解器的 benchmark 协议；
    \item 开源代码、query logs、candidate library 和实验脚本；
    \item 2--4 篇高水平论文，包括至少一篇 AI 顶会或 CCF-A 会议目标论文。
\end{enumerate}

\section{结论}

本课题的核心判断是：单纯让 LLM 直接求解组合优化问题，容易退化为黑盒启发式；单纯做某个经典优化算法的展开，又可能缺少会议新意。更有潜力的方向是将二者结合：\textbf{让 LLM 在受控、可验证、可理论分析的算法空间中，离线设计专业求解器可使用的 operator-splitting unrolled optimiser 组件。}

该方向具有三方面优势：
\begin{itemize}[leftmargin=2em]
    \item \textbf{方法新意}：连接 LLM 自动算法设计与 algorithm unrolling；
    \item \textbf{理论抓手}：通过 safe policy class 分析收敛、残差稳定和误差传播；
    \item \textbf{工程价值}：面向组合优化专业求解器中的 LP/QP 松弛、LNS、分解子问题和 repair 模块。
\end{itemize}

因此，本课题既具有博士研究的深度和持续性，也具有面向 AI 会议和专业优化社区的双重投稿潜力。

\newpage
\begin{thebibliography}{99}

\bibitem{funsearch}
B. Romera-Paredes, M. Barekatain, A. Novikov, et al.
\newblock Mathematical discoveries from program search with large language models.
\newblock \emph{Nature}, 2024.

\bibitem{eoh}
F. Liu, X. Tong, M. Yuan, X. Lin, F. Luo, Z. Wang, Z. Lu, and Q. Zhang.
\newblock Evolution of heuristics: Towards efficient automatic algorithm design using large language model.
\newblock \emph{Proceedings of ICML}, 2024.

\bibitem{reevo}
H. Ye, J. Wang, Z. Cao, F. Berto, C. Hua, H. Kim, J. Park, and G. Song.
\newblock ReEvo: Large language models as hyper-heuristics with reflective evolution.
\newblock \emph{Proceedings of NeurIPS}, 2024.

\bibitem{llmlns}
Large Language Model-driven Large Neighborhood Search for Large-Scale MILP Problems.
\newblock \emph{ICML}, 2025.

\bibitem{pdhgnet}
B. Li, L. Yang, Y. Chen, S. Wang, Q. Chen, H. Mao, Y. Ma, A. Wang, T. Ding, J. Tang, and R. Sun.
\newblock PDHG-Unrolled Learning-to-Optimize Method for Large-Scale Linear Programming.
\newblock \emph{Proceedings of ICML}, 2024.

\bibitem{deepdistributedqp}
A. D. Saravanos, H. Kuperman, A. Oshin, A. T. Abdul, V. Pacelli, and E. A. Theodorou.
\newblock Deep Distributed Optimization for Large-Scale Quadratic Programming.
\newblock \emph{ICLR}, 2025.

\bibitem{deepflexqp}
A. Oshin, R. V. Ghosh, A. D. Saravanos, and E. A. Theodorou.
\newblock Deep FlexQP: Accelerated Nonlinear Programming via Deep Unfolding.
\newblock \emph{ICLR}, 2026.

\bibitem{cobench}
W. Sun, S. Feng, S. Li, and Y. Yang.
\newblock CO-Bench: Benchmarking Language Model Agents in Algorithm Search for Combinatorial Optimization.
\newblock \emph{Proceedings of AAAI}, 2026.

\bibitem{heurigym}
H. Chen, Y. Wang, Y. Cai, H. Hu, J. Li, S. Huang, C. Deng, R. Liang, S. Kong, H. Ren, S. Samaranayake, C. Gomes, and Z. Zhang.
\newblock HeuriGym: An Agentic Benchmark for LLM-Crafted Heuristics in Combinatorial Optimization.
\newblock \emph{ICLR Benchmark Track}, 2026.

\bibitem{miplib2017}
A. Gleixner, G. Hendel, G. Gamrath, et al.
\newblock MIPLIB 2017: Data-driven compilation of the 6th mixed-integer programming library.
\newblock \emph{Mathematical Programming Computation}, 2021.

\bibitem{pdlp}
D. Applegate, M. Díaz, H. Lu, and M. Lubin.
\newblock Practical Large-Scale Linear Programming using Primal-Dual Hybrid Gradient.
\newblock \emph{NeurIPS}, 2021.

\bibitem{chambollepock}
A. Chambolle and T. Pock.
\newblock A first-order primal-dual algorithm for convex problems with applications to imaging.
\newblock \emph{Journal of Mathematical Imaging and Vision}, 2011.

\bibitem{admm}
S. Boyd, N. Parikh, E. Chu, B. Peleato, and J. Eckstein.
\newblock Distributed Optimization and Statistical Learning via the Alternating Direction Method of Multipliers.
\newblock \emph{Foundations and Trends in Machine Learning}, 2011.

\bibitem{osqp}
B. Stellato, G. Banjac, P. Goulart, A. Bemporad, and S. Boyd.
\newblock OSQP: An operator splitting solver for quadratic programs.
\newblock \emph{Mathematical Programming Computation}, 2020.

\bibitem{ml4co}
M. Gasse, D. Chételat, N. Ferroni, L. Charlin, and A. Lodi.
\newblock The Machine Learning for Combinatorial Optimization Competition.
\newblock \emph{Proceedings of Machine Learning Research}, 2022.

\end{thebibliography}

\end{document}
